\section{Agent Roles and Decision Interfaces}
\label{sec:roles}

\begin{designbox}
Every scientific and scheduling decision crosses a narrow, typed interface --- a
Python dataclass with named fields --- so that authority is explicit, verdicts are
machine-checkable, and no role can quietly reach past its contract. Four
persistent roles carry all judgment; the infrastructure only routes their typed
outputs.
\end{designbox}

Argus runs four persistent roles. Three form a numbered execution spine ---
Planner (L4), Engineer (L1), Reviewer (L2) --- while the Manager sits outside the
numbering because it spans the whole pipeline rather than occupying a station on
it. (A fifth role, the Curator, runs only in team/subagent mode and is out of
scope here.) This section documents each role's \emph{decision interface}: the
typed contract it emits and the authority it holds.

\subsection{Manager --- front door and stage authority}

The Manager converts an operator's free text into a committed campaign in a single
grounded model call that emits three structured axes at once ---
\code{CONFIG} (e.g.\ a backend/model change), \code{CONTROL} (whether to dispatch
at all), and \code{ROUTE} (chat vs.\ task and which vertical). A
\code{CONTROL: NO\_DISPATCH} decision is authoritative: the bridge forces a
read-only self-reply and, if that fails, fails closed rather than enqueuing work.
The result of a task division is a \code{Division} (task, chosen vertical, stage
template); stage progression is a \code{StageTransition}. The Manager is the sole
writer of \code{current\_stage}; the transition's \code{action} and \code{source}
fields (Table~\ref{tab:manager}) make both the decision and its justification
auditable, including the several fail-safe holds that keep the pipeline coherent
when a downstream role is unavailable.

\begin{table}[h]
\centering
\small
\begin{tabularx}{\linewidth}{@{}p{2.6cm} p{2.2cm} X@{}}
\toprule
\textbf{Field} & \textbf{Type / values} & \textbf{Meaning} \\
\midrule
\multicolumn{3}{@{}l}{\textit{Division} \; \code{manager/\_core.py}} \\
\code{task}     & str  & the committed objective text \\
\code{vertical} & str  & the one selected task shape (never re-classified) \\
\code{stages}   & list & the vertical's ordered stage template \\
\addlinespace
\multicolumn{3}{@{}l}{\textit{StageTransition} \; \code{manager/\_core.py}} \\
\code{action}   & advance $\vert$ hold $\vert$ rollback $\vert$ complete & the stage move applied to \code{current\_stage} \\
\code{source}   & manager\_llm $\vert$ *\_hold & why: a model decision, or a fail-safe hold (no review / no runner / failsafe / illegal target) \\
\bottomrule
\end{tabularx}
\caption{Manager decision interface. A \code{hold} writes nothing; only the
Manager mutates \code{current\_stage}.}
\label{tab:manager}
\end{table}

\subsection{Planner (L4) --- forward scheduling}

When the backlog empties in continuous mode, the Planner proposes the next batch
of work as a set of \code{TaskSpec}s and returns a \code{PlannerVerdict}
(Table~\ref{tab:planner}). Tasks form an optional DAG: each \code{TaskSpec} has a
local \code{key} and a \code{deps} list, which the scheduler maps to real backlog
ids when it enqueues the batch (flat tasks leave both empty). Each task carries a
structured \code{scope} --- \code{bounded} or \code{final\_submission} --- that is
threaded to the Reviewer as a backlog tag, so completion authority never has to
re-parse prose to learn a task's scope. The verdict also has a first-class,
intentional \code{waiting} outcome (with a durable \code{WaitingContract}:
a blocker fingerprint, a recheck condition, and a token) for the case where the
project is correctly blocked on a live external job and there is no genuinely new
high-impact work to queue --- an idle state that is neither an error nor make-work.
The Planner owns the per-stage checklist and emits \code{checklist\_ops} to edit
it. The canonical verdict-status contract and its (success, recoverable) policy
live in \srcpath{core/planner\_verdict.py} (Appendix~\ref{app:status}).

\begin{table}[h]
\centering
\small
\begin{tabularx}{\linewidth}{@{}p{2.7cm} p{2.3cm} X@{}}
\toprule
\textbf{Field} & \textbf{Type} & \textbf{Meaning} \\
\midrule
\multicolumn{3}{@{}l}{\textit{TaskSpec} \; \code{planner/planner.py}} \\
\code{title, objective} & str & human title + full actionable description \\
\code{impact\_score}    & int 0--5 & parser accepts only high-value work \\
\code{scope}            & bounded $\vert$ final\_submission & threaded to the Reviewer as a tag \\
\code{key, deps}        & str, list & DAG node id and its dependencies \\
\addlinespace
\multicolumn{3}{@{}l}{\textit{PlannerVerdict} \; \code{planner/planner.py}} \\
\code{project\_done}    & bool & project-level completion proposal \\
\code{new\_tasks}       & list[TaskSpec] & the next batch \\
\code{waiting}          & bool & intentional idle on a live external blocker \\
\code{waiting\_reason}  & str & paired with a \code{WaitingContract} \\
\code{restart\_daemon}  & bool & request a controlled daemon restart \\
\code{checklist\_ops}   & list & per-stage checklist edits (Planner owns it) \\
\bottomrule
\end{tabularx}
\caption{Planner decision interface (abridged). Token-usage fields are omitted;
scope is structured, never inferred from the objective text.}
\label{tab:planner}
\end{table}

\subsection{Engineer (L1) --- one round of real work}

The Engineer executes a single round against real tools and hardware. Its
interface to a provider is the per-call \code{RunnerOptions} contract, and it
receives back a \code{RunnerResult} (Table~\ref{tab:engineer}). Two properties of
this contract matter for reliability and integrity. First, provider-side spend
fences (\code{max\_budget\_usd}, \code{max\_ai\_credits}) are populated from an
atomic budget \emph{reservation}, not from role configuration, so a role cannot
raise its own ceiling (Section~\ref{sec:execution}). Second, the result carries
usage-\emph{presence} booleans so that a real zero token count is distinguishable
from a provider that simply omitted usage --- the evidence plane never invents a
number a provider did not report. Watchdog and streaming hooks
(\code{watchdog\_soft/hard\_idle\_seconds}, \code{on\_agent\_message}) are honored
by the subprocess backends and ignored by the deterministic test backend. The
Engineer's own round-loop configuration (round bounds, stall thresholds, effective-progress
timeout) is documented with the reliability guards in Section~\ref{sec:reliability}.

\begin{table}[h]
\centering
\small
\begin{tabularx}{\linewidth}{@{}p{3.3cm} X@{}}
\toprule
\textbf{Field} & \textbf{Meaning} \\
\midrule
\multicolumn{2}{@{}l}{\textit{RunnerOptions} (per-call backend knobs) \; \code{core/models.py}} \\
\code{model, reasoning\_effort} & provider model and effort level \\
\code{output\_schema\_path}     & structured-output schema for the call \\
\code{live\_search}             & provider web search; research stage only \\
\code{sandbox\_mode}            & containment policy for the spawned role \\
\code{max\_budget\_usd, max\_ai\_credits} & spend fences from a reservation, not config \\
\code{watchdog\_*\_idle\_seconds, on\_agent\_message} & idle watchdog + streaming (subprocess backends) \\
\addlinespace
\multicolumn{2}{@{}l}{\textit{RunnerResult} (per-call outcome) \; \code{core/models.py}} \\
\code{call\_id, call\_id\_log\_correlated} & audit correlation to the event tape \\
\code{*\_tokens\_present}       & presence flags: a real zero vs.\ an omitted usage \\
\code{premium\_requests, total\_nano\_aiu, cost\_usd} & usage and cost accounting \\
\code{pricing\_status, tool\_activity\_observed} & pricing state; whether the round used tools \\
\bottomrule
\end{tabularx}
\caption{Engineer per-call backend contract (abridged). Spend fences come from a
reservation; presence flags keep usage accounting honest.}
\label{tab:engineer}
\end{table}

\subsection{Reviewer (L2) --- the completion authority}

The Reviewer is the sole authority on completion, and there is no separate
hardcoded completion gate anywhere in the system. Each round it inspects the
Engineer's output against a structured stage checklist and returns a
\code{ReviewDecision} (Table~\ref{tab:reviewer}) whose \code{status} is
\code{done}, \code{continue}, or \code{blocked}. Because the Reviewer runs in the
same work-tree as the Engineer and is given the mission's execution log with grep
recipes, it can audit \emph{how} a result was produced --- catching a hardcoded
answer, a skipped step, or a faked metric --- rather than trusting a summary
(Section~\ref{sec:evidence}). The verdict is rich: \code{progress\_class} is the
Reviewer's structured judgment of forward motion (the infrastructure counts it,
never infers it); \code{operator\_question} is the single human-facing block
question; \code{scope} plus \code{checklist} form the final-submission gate;
\code{skill\_ops} and \code{wiki\_ops} are authored here with no Manager gate; and
\code{backend\_unavailable} flags an infrastructure failure so it is never mistaken
for a \code{continue}. A \code{final\_submission\_certified} property is
fail-closed --- it is true only when the structured gate is genuinely satisfied.

\begin{table}[h]
\centering
\small
\begin{tabularx}{\linewidth}{@{}p{3.4cm} X@{}}
\toprule
\textbf{Field} & \textbf{Meaning} \\
\midrule
\code{status} & \code{done} $\vert$ \code{continue} $\vert$ \code{blocked} (the completion verdict) \\
\code{progress\_class} & \code{decision} $\vert$ \code{evidence} $\vert$ \code{setup\_only} $\vert$ \code{artifact\_sync\_only} $\vert$ \code{none} \\
\code{next\_action} & concrete next step handed to the subsequent Engineer round \\
\code{operator\_question} & the one human-facing block question, if any \\
\code{scope, checklist} & final-submission gate (fail-closed) \\
\code{achievement, failure\_cause} & certified achievement payload / diagnosed cause \\
\code{skill\_ops, wiki\_ops} & knowledge-system edits (Reviewer is sole author) \\
\code{backend\_unavailable} & infrastructure-failure flag, not a judgment \\
\code{step\_back} & anti-plan-lock-in signal \\
\bottomrule
\end{tabularx}
\caption{Reviewer verdict --- the completion contract (\code{core/models.py},
\code{reviewer/reviewer\_schema.json}). All fields are required by the schema; the
achievement payload is emitted only on a \code{done} verdict.}
\label{tab:reviewer}
\end{table}

\begin{callout}[Why ``done'' has exactly one owner]
Premature acceptance (Section~\ref{sec:episodic}) is answered structurally:
completion is a single typed verdict from a role that sees the artifacts and the
execution log, not a self-report from the role that did the work, and not a
keyword the infrastructure fished out of a summary. If the Reviewer is wrong, that
is a model-quality limitation (Section~\ref{sec:limitations}) --- but the authority
is never ambiguous, and it is never silently overridden.
\end{callout}

\subsection{What was deliberately removed}

The commitment to keeping judgment with the agent is visible as deletions, not
slogans. A prior reviewer path
scanned the Engineer's summary with regexes and forcibly rewrote a \code{done}
verdict into \code{continue}; it and its pattern constants were removed, and the
underlying constraint now lives in the Reviewer's own prompt. Chat-vs-task routing
is a low-effort model call, not a length or regex heuristic. And task type,
lifetime, and scope are decided by explicit structured signals --- a vertical's
completion gate, a configuration flag, a backlog tag --- not by grepping the
objective text.
